\chapter*{Conclusion générale}
\addcontentsline{toc}{chapter}{Conclusion générale}

Le stage réalisé au sein de \CompanyName{} avait pour objectif de concevoir et
développer un outil de visualisation et de suivi de la disponibilité des bornes
de recharge électrique. Ce sujet a conduit à dépasser la simple représentation
d'un état enregistré en base : une disponibilité exploitable doit être dérivée
des messages de la borne, de la fraîcheur de sa communication, de ses
connecteurs, de ses transactions et des contraintes d'exploitation.

\section*{Bilan du travail réalisé}

\ProjectName{} aboutit à un MVP vertical qui relie une borne simulée à une
application métier complète. Le simulateur SAP communique en OCPP~1.6J avec une
passerelle Python ; les événements sont transmis de manière signée à Laravel,
persistés puis projetés en états de disponibilité. Ces états alimentent la
cartographie, les alertes, les tableaux de bord et le parcours de recharge.
Une session peut être préparée par le client, démarrée à distance, suivie à
partir des \texttt{MeterValues}, arrêtée puis facturée au moyen d'un prestataire
HTTP simulé derrière l'interface \texttt{PaymentGateway}.

Le produit couvre également l'identité locale et Google, les invitations, les
demandes de démonstration, l'isolation multi-organisation, les stations et
connecteurs, les commandes OCPP, les alertes, les interventions, la maintenance,
les tarifs, les abonnements clients, le contrat commercial des organisations et
les rapports. Les espaces Super Administrateur, Administrateur, Opérateur,
Technicien et Client présentent des indicateurs et des actions adaptés à leurs
responsabilités, au lieu de réutiliser un tableau de bord générique.

La validation locale confirme la cohérence du socle réalisé : 203 tests backend
et 1\,980 assertions réussissent, la passerelle et les outils OCPP disposent de
leurs propres tests, le lint frontend est valide et la compilation de production
aboutit.
Ces résultats ne remplacent pas une recette sur infrastructure réelle, mais ils
fournissent une base mesurable et reproductible pour le MVP.

\section*{Apports du stage}

Ce travail a permis de mobiliser les différentes étapes d'un projet de génie
logiciel : clarification des acteurs et des règles métier, construction d'un
Product Backlog, modélisation UML, choix d'une architecture, développement
organisé en Sprints, intégration de services externes, tests et documentation. Il a
également renforcé des compétences techniques en React et TypeScript, Laravel,
PostgreSQL, WebSocket, OCPP, Docker, sécurité multi-tenant et génération de
documents.

L'apport principal ne se limite pas au volume fonctionnel. Le projet met en place
des frontières explicites entre l'interface, le domaine et l'infrastructure :
Laravel reste propriétaire des autorisations et des règles, la passerelle OCPP
ne calcule pas l'état métier, le frontend n'invente pas les mesures et le
paiement dépend d'un contrat remplaçable. Cette séparation réduit le risque lors
du remplacement futur d'un simulateur par un service ou un équipement réel.

\section*{Difficultés et réponses apportées}

Les principales difficultés ont concerné les dépendances entre les objectifs des Sprints.
Une session de recharge fiable dépend du commissioning, de la connexion OCPP,
des états de connecteur, du tarif et du paiement. La progression par Sprints a donc
été organisée sans supposer que les fonctionnalités étaient indépendantes :
chaque Sprint consolidait les contrats nécessaires au suivant.

La deuxième difficulté a porté sur la distinction entre données simulées et
preuves réelles. L'interface n'affiche pas de télémétrie synthétique lorsque la
borne ne fournit pas de \texttt{MeterValues}. De même, un bouton de commande
distante ne signifie pas que l'action a réussi ; le résultat provient de la
réponse OCPP. Enfin, WireMock simule le comportement d'une API financière, mais
il n'est jamais présenté comme un paiement bancaire certifié.

L'isolation des organisations, les erreurs asynchrones et les différences entre
SQLite de test et PostgreSQL ont également nécessité des contrôles spécifiques.
Les Form Requests, policies, services de périmètre, clés d'idempotence, jobs
planifiés et tests de non-régression ont été utilisés pour traiter ces risques.

\section*{Limites de la version actuelle}

La version présentée reste un MVP exécuté dans un environnement local. La borne
est simulée, le paiement est externe mais fictif et le déploiement de production
n'est pas réalisé. Les mesures de performance actuelles ne constituent pas une
campagne de charge représentative d'un grand réseau. La couverture fonctionnelle
est importante, mais une recette finale avec l'encadrant, des tests E2E plus
larges et l'observabilité de production restent nécessaires avant une
exploitation commerciale. La spécification OpenAPI, les manuels et le cahier de
tests sont désormais versionnés avec le code.

\section*{Perspectives}

Les évolutions doivent être introduites selon leur dépendance et leur niveau de
risque. La première priorité est l'industrialisation : CI/CD, TLS, gestion
centralisée des secrets, sauvegardes, métriques, journaux corrélés, alertes
techniques et tests de charge OCPP. Une recette formelle pourra ensuite figer les
critères d'acceptation du MVP.

La deuxième priorité consiste à confronter les adaptateurs à des systèmes réels :
connecter une borne physique OCPP~1.6J et intégrer, dans un environnement de test,
un prestataire compatible avec le marché tunisien. Cette étape devra couvrir la
réconciliation, les remboursements, les litiges et les exigences contractuelles
du fournisseur.

À moyen terme, la plateforme pourra enrichir la cartographie et les itinéraires,
étudier OCPP~2.0.1 et préparer une gestion de firmware signée avec suivi de
progression et retour arrière. Les fonctions d'intelligence artificielle,
notamment la prédiction de pannes ou l'assistance au diagnostic, ne devront être
ajoutées qu'après collecte de données fiables, définition d'indicateurs
mesurables et stabilisation du cœur opérationnel.

Ainsi, \ProjectName{} fournit une base cohérente pour superviser un réseau de
recharge, coordonner les acteurs et faire évoluer les intégrations sans remettre
en cause les règles essentielles du domaine. Le stage a permis de transformer un
sujet initial de visualisation en une plateforme démontrable, testée et ouverte
à une validation progressive sur des infrastructures réelles.
